\documentclass[12pt]{article}
\usepackage{amsmath, amssymb, amsthm}
\usepackage{hyperref}

\title{BSD Certificate for $J_0(143)$:\\
  H4 Invariant Labels and M* Normalization}
\author{Opera Numerorum Archive}
\date{2026}

\newtheorem{theorem}{Theorem}
\newtheorem{lemma}[theorem]{Lemma}
\newtheorem{definition}[theorem]{Definition}
\newtheorem{remark}[theorem]{Remark}

\begin{document}
\maketitle

\begin{abstract}
This certificate records the three distinct H4-related quantities
that appear in the BSD verification for $J_0(143)$, and insists on
unambiguous labels for each.  The M*~normalization ratio $\approx 12/11$,
the H4 base $120/11$, and the H4 Coxeter-group order $14400$ are
mathematically independent; conflating their labels corrupts the
argument.
\end{abstract}

\section{Three Distinct H4 Quantities}

The BSD certificate for $J_0(143)$ involves three quantities that
each carry the prefix ``H4'' in informal usage.  They must not be
merged.

\begin{definition}[M* normalization ratio]\label{def:mstar-ratio}
The \emph{M*~normalization ratio} is the theoretical target
\[
  r_{M^*} \;=\; \frac{12}{11} \;\approx\; 1.0909\ldots
\]
It is the theoretical value that the computed ratio
$M^*_{\mathrm{raw}}/B_{H4}$ is claimed to approximate across T-22
sequences.  Measured values of $M^*_{\mathrm{raw}}$ range from
approximately $11.976$ to $12.003$, giving
\[
  \frac{M^*_{\mathrm{raw}}}{B_{H4}}
  \;=\;
  \frac{11.98\text{ (approx)}}{120/11}
  \;=\;
  \frac{11.98 \times 11}{120}
  \;\approx\; 1.098.
\]
This lies within roughly $0.8\text{--}1\%$ of $12/11\approx 1.0909$;
the two are numerically close but not equal.

\textbf{Caution:} $M^*_{\mathrm{raw}}$ is \emph{not} exactly $12$.
If it were, the exact quotient would be $12/(120/11)=11/10=1.1$, not
$12/11\approx 1.0909$.  The label $r_{M^*}=12/11$ records a
theoretical target, not a consequence of an exact algebraic identity.

This quantity is \emph{not} the H4 base and \emph{not} the group order.
\end{definition}

\begin{definition}[H4 base]\label{def:h4-base}
The \emph{H4 base} is
\[
  B_{H4} \;=\; \frac{120}{11} \;\approx\; 10.9091\ldots
\]
It is the normalization divisor used to form the dimensionless ratio
$r_{M^*} \approx 12/11$.  The label \texttt{H4\_base} is reserved
exclusively for $120/11$.  It is \emph{not} the M*~ratio $12/11$ and
\emph{not} the group order $14400$.
\end{definition}

\begin{definition}[H4 Coxeter-group order]\label{def:h4-order}
The \emph{H4 Coxeter-group order} is
\[
  |W(H4)| \;=\; 14400.
\]
This is the order of the finite reflection group $W(H4)$ acting on
$\mathbb{R}^4$.  It satisfies the layer-ratio identity
\[
  \frac{|W(H4)|}{|W(H3)|} \;=\; \frac{14400}{120} \;=\; 120,
\]
which equals the number of M8G layers.  This quantity is \emph{not}
the ratio $12/11$ and \emph{not} the base $120/11$.
\end{definition}

\section{Label Integrity for Archive Regeneration}

\begin{remark}[Forbidden shorthands]\label{rem:forbidden}
A bare, unqualified label that assigns the value $12/11$ to the key
\texttt{H4}---without a qualifier suffix---is forbidden in this archive
because it is ambiguous: it could denote the (approximate) M*~ratio,
a mis-stated H4 base, or a shorthand for the group order.  Any
regenerated certificate must use:
\begin{itemize}
  \item \texttt{Mstar\_ratio} for $r_{M^*} \approx 12/11$,
  \item \texttt{H4\_base} for $B_{H4} = 120/11$ exactly,
  \item \texttt{H4\_coxeter\_group\_order} for $|W(H4)| = 14400$.
\end{itemize}
A bare denominator label (e.g.\ ``$11$'' standing alone as an H4
label) is likewise rejected; the full fraction $120/11$ must appear.
\end{remark}

\section{BSD Anchor: $\Delta_{DS}^{(4)} / B_{H4}$}

The key BSD match for $J_0(143)$ is

\[
  \frac{\Delta_{DS}^{(4)}}{B_{H4}}
  \;=\; 2.1812
  \;\approx\; 2 \cdot r_{M^*}
  \;=\; 2 \cdot \frac{12}{11}
  \;=\; \frac{24}{11}
  \;\approx\; 2.1818,
  \qquad \text{error } 0.027\%.
\]

Here the \emph{M*~normalization} $r_{M^*} \approx 12/11$ appears
explicitly as a factor on the right-hand side, separated from the
H4 base $B_{H4} = 120/11$ that appears as the divisor on the left.
The two roles are structurally distinct and must be labeled as such in
any automated certificate.

\section{Coxeter Unification Summary}

The H4 Coxeter group unifies five certified constants:
\begin{center}
\begin{tabular}{lll}
\hline
Constant & Value & Source \\
\hline
$r_{M^*}$ (M*~ratio, approx.) & $\approx 12/11$ & M21, M22, M23 \\
$B_{H4}$ (H4 base, exact)     & $120/11 \approx 10.9091$ & M21, M22, M23\_BSD \\
$|W(H4)|$ (group order, exact) & $14400$ & M8G, Section~8 \\
$\Omega/R$                     & $11.929 \approx 12$ (err $0.59\%$) & M23 \\
$h_2$                          & $10$ & M6 \\
\hline
\end{tabular}
\end{center}

\noindent
None of these five entries is interchangeable with another.  The
archive check script \texttt{pistus-theoria/check\_h4\_labels.py}
enforces these distinctions automatically.

\end{document}
